% Generated draft: requires mathematical and source review.
\documentclass[11pt]{article}
\usepackage[margin=25mm]{geometry}
\usepackage{fontspec}
\setmainfont{Latin Modern Roman}
\usepackage{amsmath,amssymb,mathtools,unicode-math}
\setmathfont{Latin Modern Math}
\usepackage{enumitem,needspace,xurl,hyperref}
\setlength{\parindent}{0pt}
\setlength{\parskip}{5pt}
\setlength{\emergencystretch}{3em}
\newcommand{\N}{\mathbb{N}}
\newcommand{\Z}{\mathbb{Z}}
\newcommand{\Q}{\mathbb{Q}}
\newcommand{\R}{\mathbb{R}}
\newcommand{\C}{\mathbb{C}}
\newcommand{\F}{\mathbb{F}}
\newcommand{\A}{\mathbb{A}}
\newcommand{\PP}{\mathbb{P}}
\newcommand{\E}{\mathbb{E}}
\newcommand{\eps}{\varepsilon}
\newcommand{\dd}{\,\mathrm{d}}
\begin{document}
\section*{UnsolvedMath ID: 800464}
\textbf{Source problem code:} catalogue entry
\subsection{Definitions and mathematical statement}
Write \((X,d,\nu)\) for the intended metric measure space and define the open metric ball by
\[
B_r(x)=\{y\in X:d(x,y)<r\},\qquad r>0.
\]
The supplied extract does not specify the hypotheses on \(X\) or \(\nu\), the meanings of the parameters \(n\) and \(\lambda\), or the intended range of radii. Ratios below require finite, strictly positive denominator measures. These requirements explain the notation; they do not recover the missing source hypotheses.

Let \(\gamma:(-2,2)\to X\) be a minimizing geodesic, meaning that its subsegments realize the distance between their endpoints. The extract does not specify its speed or normalization. Conjecture 8, attributed in the supplied record to Naber (2020), gives the following two assertions for \(s,t\in(-1,1)\):
\begin{enumerate}
\item \textbf{Lipschitz rate, as transcribed.} For a positive radius \(r\), whose intended range and quantification are not supplied,
\[
\frac{\nu(B_r(\gamma(s)))}{\nu(B_r(\gamma(t)))}
\leq C(n,\lambda)|t-s|.
\]
The constant \(C(n,\lambda)\) is not further specified. This formula cannot hold as written whenever the ratio at \(s=t\) is defined and positive: it then reads \(1\leq0\). No exclusion of \(s=t\) is supplied. Thus the displayed formula is retained as source text requiring correction or clarification, rather than asserted as a consistent conjecture. The intended Lipschitz estimate cannot be uniquely recovered from this extract.
\item \textbf{Constant density, as transcribed.} If \(X\) is noncollapsed, then for every \(s,t\in(-1,1)\),
\[
\lim_{r\downarrow0}
\frac{\nu(B_r(\gamma(s)))}{\nu(B_r(\gamma(t)))}=1.
\]
Here the ratios must be defined for sufficiently small positive radii. The intended definition of noncollapsed and the ambient hypotheses needed for this assertion are absent from the extract.
\end{enumerate}
Consequently, the supplied record does not determine a complete, consistent mathematical formulation of both assertions.
\subsection{Short English statement}
The supplied conjecture concerns ratios of measures of small balls centered along the interior of a minimizing geodesic. It states a Lipschitz estimate and, under an unspecified noncollapse assumption, convergence of these ratios to \(1\). The first displayed estimate is inconsistent at equal parameters, and the ambient hypotheses are missing; the intended corrected statement requires source review.
\subsection{Sources}
\begin{itemize}
\item Local UnsolvedMath registry, permanent ID \(800464\), source problem code \texttt{catalogue entry}, in the AMR Open Problem Lists collection. The supplied record labels the problem \texttt{open}; this is a supplied status label, not an independently verified mathematical or literature assessment.
\item Naber, \emph{Conjectures and Open Questions on Spaces with Lower Ricci Curvature} (2020), Conjecture 8, as identified by the supplied record: \url{https://arxiv.org/abs/2010.10031}. Attribution and proposed year are supplied by the registry. The background records access on 2026-10-16 and extraction method \texttt{arxiv\_latex}; these are supplied provenance statements, not a claim of independent access or verification.
\item Supplied upstream revision: \texttt{d371ee917ac9647fbb06cd6fc529deec6f3dda82}.
\end{itemize}
\end{document}
